\documentclass[11pt]{article}
\usepackage[margin=0.3in,top=0.25in,bottom=0.4in]{geometry}
\makeatletter
\renewcommand{\@maketitle}{\newpage\null\vskip -3.5em\begin{center}{\LARGE \@title \par}\end{center}\par\vskip 0.2em}
\makeatother
\usepackage{amsmath,amssymb}
\usepackage{hyperref}

\title{What \texttt{flux2d} computes: photon-flux-convolved exclusive and\\
diffractive $q\bar q$(+g) production, differential in $p_T$ and $y$}
\author{}
\date{}

\begin{document}
\maketitle

\section{Overview}

\texttt{flux2d} computes, on a genuine 2D grid of the produced charm
quark's transverse momentum $p_T$ and rapidity $y$, two photon-nucleus
ultraperipheral collision (UPC) cross sections:

\begin{itemize}
\item \textbf{Exclusive}: $\gamma^*A\to q\bar qA$, a $q\bar q$ dijet with
no radiated gluon. Function: \texttt{exclusiveCrossSection()}
(\texttt{src/int\_exclusive.cpp}).
\item \textbf{Diffractive}: $\gamma^*A\to q\bar q g\,A$, the same process
with an additional radiated gluon of transverse momentum $k$, carrying
Pomeron momentum fraction $x_{\mathbb P}$. Function:
\texttt{diffractiveCrossSection\_xpo()}
(\texttt{src/int\_diffractive\_xpo.cpp}).
\end{itemize}

\textbf{Note the two different momenta}: $p_T$ is the produced charm
quark's own transverse momentum — a fixed input at every grid point, the
same physical quantity in both processes. $k$ is the diffractive process's
\emph{extra radiated gluon}'s transverse momentum — it exists only in the
diffractive process, is integrated out (not observed), and is capped at
$p_T$ (Section~\ref{sec:diffractive}). The two are never the same variable
under two names.

In both cases the equivalent photon flux from the nucleus is
\emph{genuinely convolved} — the photon energy $qp$ and impact parameter
$b$ are real Monte Carlo (VEGAS) integration variables, weighted by
\texttt{photon\_flux()} (\texttt{src/gamma\_aa.cpp}) — as opposed to the
fixed-$W$ approximation used elsewhere in this project family, which
assigns the photon a single analytic energy and drops the flux integral
entirely.

Both results are differential in exactly $d^2p_T\,dy$: $p_T$ and $y$ are
fixed inputs at each grid point, not integrated over. This is what makes
the two processes directly comparable, and is exactly the variable set of
the original \texttt{main\_compare.cpp} comparison in \texttt{exclusive/}.

\section{Common kinematics}

At each grid point $(p_T,y)$:
\[
m_T = \sqrt{p_T^2+m^2}, \qquad p^+ \equiv pp = \frac{m_T}{\sqrt2}\,e^{y},
\]
where $m$ is the heavy-quark mass (charm, $m=1.5$~GeV). The photon energy
$qp\equiv q^+$ is sampled by VEGAS, uniformly in $[pp,\,qp_{\max}]$:
\[
qp = pp + u_{qp}\,(qp_{\max}-pp), \qquad u_{qp}\in[0,1],
\]
with Jacobian $\partial qp/\partial u_{qp} = qp_{\max}-pp$. The light-cone
momentum fractions of the quark and antiquark are
\[
z_1 = \frac{pp}{qp}, \qquad z_2 = 1-z_1,
\]
and the photon-nucleus invariant mass and $q\bar q$ invariant mass are
\[
W^2 = \sqrt2\,qp\,ss, \qquad M_{q\bar q}^2 = \frac{p_T^2+m^2}{z_1 z_2},
\]
where $ss=\sqrt{s}=5360$~GeV is the fixed photon-nucleus c.o.m.\ energy.
The impact parameter $b$ is sampled in $[b_{\min},b_{\max}]$ and enters
only through the photon flux (Section~\ref{sec:flux}).

\section{Exclusive process}

$x_{\mathbb P}$ is not an independent variable here — with no radiated
gluon, the whole Pomeron momentum goes into the $q\bar q$ pair, so
\[
x_{\mathbb P} = \frac{M_{q\bar q}^2}{W^2},
\]
fixed once $p_T,y,qp$ are. The integrand enforces the standard small-$x$/
dipole-applicability window $0<x_{\mathbb P}\le0.1$ (values above are set
to zero — the BK/dipole picture is not assumed to apply there), and freezes
the dipole amplitude at its initial-condition value above
$x_{\mathbb P}=0.01$ (\texttt{xP\_dip = min(xP, 0.01)}), since the
evolution table does not extend above that.

The hard kernel involves the dipole amplitude $N(r,x_{\mathbb P})$ at two
transverse dipole sizes $r_1,r_2$ (also VEGAS-integrated, in $[0,99]$
GeV$^{-1}$) and two radial Bessel kernels — a spin non-flip+flip piece
($\nu=1$) and an ``aligned-jet''/mass piece ($\nu=0$):
\[
f_{1}(r) = r\,J_1(p_Tr)\,K_1(mr)\,N(r,x_{\mathbb P}), \qquad
f_{0}(r) = r\,J_0(p_Tr)\,K_0(mr)\,N(r,x_{\mathbb P}),
\]
combined as
\[
g = (z_1^2+z_2^2)\,f_1(r_1)f_1(r_2) + f_0(r_1)f_0(r_2).
\]
The full integrand (\texttt{integrand\_exclusive()},
\texttt{src/integrand.cpp}) is
\[
\frac{d\sigma^{\rm excl}}{d^2p_T\,dy} \;\propto\;
\int_0^{99}\!\!dr_1\int_0^{99}\!\!dr_2\int_{pp}^{qp_{\max}}\!\!dqp
\int_{b_{\min}}^{b_{\max}}\!\!db\;\;
z_1\,m^2\,\big(qp_{\max}-pp\big)\,\text{photon\_flux}(b,qp)\,g,
\]
a 4D VEGAS integral over $\{r_1,r_2,u_{qp},b\}$
(\texttt{exclusiveCrossSection()}). \textbf{This result is missing the
overall prefactor} $\alpha_{\rm em}\,N_c\,e_f^2\,S_\perp/(2\pi^2)$, applied
downstream (e.g.\ in the plotting scripts) rather than in the C++ code.

\section{Diffractive process}
\label{sec:diffractive}

Here the Pomeron's momentum fraction $x_{\mathbb P}$ is an
\emph{independent} physical variable (the momentum carried by the radiated
gluon plus the $q\bar q$ system), distinct from the internal struck-parton
variable
\[
x = \frac{M_{q\bar q}^2}{x_{\mathbb P}\,W^2}, \qquad 0<x<1.
\]
$y$ is held \textbf{fixed} at the grid value (like the exclusive case —
\texttt{par->pp} is set directly from the grid's $p_T,y$ before the VEGAS
call), and $x_{\mathbb P}$ is instead \textbf{genuinely integrated}
(summed), log-uniformly over $[10^{-6},\,0.1]$:
\[
x_{\mathbb P} = x_{\mathbb P,\min}\,e^{u_{x_{\mathbb P}}\ln(0.1/x_{\mathbb P,\min})},
\qquad u_{x_{\mathbb P}}\in[0,1].
\]
(The \texttt{exclusive/} codebase this was copied from also has a
\texttt{diffractiveCrossSection\_dlnxpo()} that does the opposite — fixes
$x_{\mathbb P}$ and integrates $y$ instead — plus fixed-$W$/no-flux
variants of both processes and an $x_{\mathbb P}$-differential exclusive
function. None of those are part of what \texttt{flux2d} computes, so
they (and the now-unused \texttt{radialIntegral()} helper they relied on)
were not copied over / were removed from this repository's source.)

The extra radiated gluon's transverse momentum $k$ — \textbf{a different
momentum from $p_T$} — is also integrated, log-mapped over
$[10^{-3},\,p_T]$~GeV. The upper bound is capped at $p_T$, the produced
charm quark's own transverse momentum, since the radiated gluon cannot
carry more transverse momentum than the produced quark itself (matching
\texttt{partonCS/src/main.cpp}'s identical \texttt{k\_upper = pt} choice
for its ``theta'' cross section):
\[
k = k_{\rm lo}\,e^{u_k\ln(p_T/k_{\rm lo})}, \qquad u_k\in[0,1].
\]
With $\omega^2 \equiv \frac{x}{1-x}k^2$, the hard factor and the
generalized dipole TMD-like piece are
\[
H^2 = \frac{z_1}{(p_T^2+m^2)^4}\Big[(z_1^2+z_2^2)\big(p_T^4+m^4\big)
+ 2m^2p_T^2\Big],
\]
\[
g_{\rm DTMD}^2 = \omega^4\,r_1 r_2\;
J_2(kr_1)K_2(\sqrt{\tfrac{x}{1-x}}\,k\,r_1)\;
J_2(kr_2)K_2(\sqrt{\tfrac{x}{1-x}}\,k\,r_2),
\]
and the full integrand (\texttt{integrand\_diffractive\_xpo()},
\texttt{src/integrand.cpp}) is
\[
\frac{d\sigma^{\rm diff}}{d^2p_T\,dy} \;=\;
\int_0^{99}\!\!dr_1 dr_2\int_{k_{\rm lo}}^{p_T}\!\!dk
\int_{pp}^{qp_{\max}}\!\!dqp \int_{b_{\min}}^{b_{\max}}\!\!db
\int_{10^{-6}}^{0.1}\!\!dx_{\mathbb P}\;\;
2\pi k\,\big(qp_{\max}-pp\big)\,\text{photon\_flux}(b,qp)\,
D_1D_2\,H^2\,\frac{g_{\rm DTMD}^2}{1-x},
\]
a 6D VEGAS integral over $\{r_1,r_2,u_k,u_{qp},b,u_{x_{\mathbb P}}\}$
(\texttt{diffractiveCrossSection\_xpo()}), with $D_1=N(r_1,\min(x_{\mathbb
P},0.01))$, $D_2=N(r_2,\min(x_{\mathbb P},0.01))$. Unlike the exclusive
side, \textbf{this result is the full physical cross section} — no
missing prefactor.

\section{Photon flux}
\label{sec:flux}

\texttt{photon\_flux(b,qp)} (\texttt{src/gamma\_aa.cpp}) implements the
equivalent-photon (Weizs\"acker--Williams) flux of a nucleus of charge $Z$,
folded with UPC survival/veto factors:
\[
\text{flux\_density}(qp,b) = \frac{\alpha_{\rm em} Z^2}{\pi^2\,\omega}\,
\frac{\eta^2}{b^2}\,K_1(\eta)^2, \qquad
\omega = \frac{qp}{\sqrt2}, \qquad
\eta = \frac{2\omega}{ss}\,m_n\,b,
\]
with $m_n=0.938$~GeV the nucleon mass. The full flux additionally
multiplies by:
\begin{itemize}
\item $\Gamma_{AA}(b)$ — a tabulated nuclear survival probability
(\texttt{data/Gamma\_AA.dat}, spline-interpolated by the \texttt{TASpline}/
\texttt{GammaAA} classes), applied for $b<150$~GeV$^{-1}$ (else $1$).
\item An electromagnetic-dissociation (EMD) selection factor depending on
the analysis channel (\texttt{par->channel = "An0n"} here): the
probability of \emph{no} mutual EM dissociation,
$P_{\rm no\,emd}=\exp(-S/b^2)$, with $S=(17.4/0.197327)^2$~GeV$^{-2}$ fixed
by the neutron-tagging requirement of the ``An0n'' channel (a single
forward neutron detected on one side, none on the other, used to tag
the photon-emitting nucleus).
\end{itemize}
This exact flux setup (nuclear/EM parameters, channel choice) matches
\texttt{exclusive/main\_flux.cpp}'s Pb-like UPC configuration at the same
$\sqrt s$.

\section{The $(p_T,y)$ grid}

\texttt{src/main\_grid.cpp} evaluates both cross sections above at every
point of a 2D grid:
\begin{itemize}
\item $p_T$: $n_{p_T}$ points, log-spaced in $[0.2,12]$~GeV.
\item $y$: $n_y$ points, linearly spaced in $[y_{\min},y_{\max}]$
(default $[0,4]$).
\end{itemize}
The grid is flattened and parallelized over both dimensions at once via
OpenMP (\texttt{\#pragma omp parallel for collapse(2)}); each $(p_T,y)$
point is an independent VEGAS call. CLI:
\[
\texttt{./flux\_grid <n\_pt> <n\_y> <calls> <y\_min> <y\_max>}
\]
(all optional, defaulting to \texttt{40 5 1e5 0 4}). Output is two plain-text
matrix files, \texttt{files/flux\_grid\_exclusive.dat} and
\texttt{files/flux\_grid\_diffractive.dat}: rows are $p_T$, columns are
$y$, with the $y$ grid values given in a
\texttt{\# pt \textbackslash y ...} header line.

\section{Known numerical caveat}

Near the exclusive side's $x_{\mathbb P}>0.1$ cutoff (Section 3), and
generally at large $p_T$, the region of $\{qp,b\}$ contributing nonzero
integrand shrinks to a thin sliver of the full VEGAS volume. At the
default \texttt{calls=1e5} this produces visible shot noise in the
high-$p_T$ tail — sharp near-vertical drops and noisy bounce-backs,
occasionally landing slightly negative (statistically consistent with
zero, not a real negative cross section). This is the same phenomenon
documented for \texttt{exclusive/main\_flux.cpp}
(\texttt{exclusive/notes/exclusive\_xpom\_application.tex}, Section~9);
raising \texttt{calls} reduces it roughly as $1/\sqrt{\texttt{calls}}$ but
does not eliminate it entirely right at the cutoff. Capping the
diffractive $k$ integral at $p_T$ (Section~\ref{sec:diffractive}) removes
this noise from the diffractive side in practice (it shrinks the VEGAS
volume exactly where the noise used to come from), but the exclusive side
is unaffected by that change and still shows it.

\section{Reproducibility}

\begin{itemize}
\item \texttt{run\_local.sh <n\_pt> <n\_y> <calls> <y\_min> <y\_max>}:
builds \texttt{flux\_grid} and runs it once.
\item \texttt{scripts/plot\_flux\_grid.py}: reads the two grid files and
plots $d\sigma/dp_T$ (both processes, one curve per $y$) against $p_T$,
with the same coupling/color-factor prefactor convention and $2\pi p_T$
Jacobian as \texttt{exclusive/scripts/plot\_flux\_spectrum.py}.
\end{itemize}

\end{document}
